%% kapa: exam=1 round=37 year=2026 period=2 form=A subject=appraisal_law subject_ko=감정평가관계법규 seq=1 number=1
%% source: Q-Net 감정평가사 기출문제 (cst00309) artlSeq=5258924 file=raw/1008/5258924/2026년 제37회 감정평가사 1차 2교시.pdf
%% conversion: kordoc 4.24.51 -> tools/gen1.py (auto); review=proofread-sonnet; flags=none
\begin{question}{1}
국토의 계획 및 이용에 관한 법령상 도시ㆍ군계획시설입체복합구역(이하 ‘입체복합구역’)에 관한 설명으로 옳은 것은? (단, 조례는 고려하지 않음)
\begin{choices}
  \item 도시ㆍ군관리계획 결정권자는 도시ㆍ군계획시설 준공 후 5년이 경과한 경우로서 해당 시설의 정비가 필요한 경우 도시ㆍ군계획시설이 결정된 토지의 전부 또는 일부를 입체복합구역으로 지정하여야 한다.
  \item 자연환경보전지역의 경우 보전관리지역에서 허용되는 범위에서 입체복합구역에서의 시설의 용도 제한을 따로 정할 수 있다.
  \item 입체복합구역에서의 건축물의 높이는 ｢건축법｣에 따른 채광 등의 확보를 위한 건축물의 높이 제한의 250퍼센트 이하의 범위에서 따로 정할 수 있다.
  \item 입체복합구역의 지정을 제안하려는 자는 대상 토지면적의 3분의 2 이상 토지소유자의 동의를 받아야 한다.
  \item 입체복합구역이 상업지역에 위치하는 경우 공간재구조화계획 입안권자는 공간재구조화계획의 입안을 위한 기초조사를 생략할 수 있다.
\end{choices}
\end{question}
\answer{5}
